%% toc: P15 165 165.10
\section{Choice rules for a coherent lower prevision}
\label{sec:annex-p15-w02b-p15-036}

Throughout, \cref{def:p15-w02b-p15-036-standing} is in force. Since
\(\underline P\) is coherent on \(B\), it equals its natural extension, so
by \cref{thm:p15-ip-natural-extension} and \cref{thm:p15-ip-dual} it has
the following properties for all \(f,g\in B\) and real constants \(c\):
\begin{enumerate}
\item[(P1)] \(\underline P(f+g)\geq\underline P(f)+\underline P(g)\);
\item[(P2)] \(\underline P(c)=c\) and \(\underline P(f+c)=\underline P(f)+c\);
\item[(P3)] \(\underline P(f)=\min_{L\in\mathcal M}L(f)\) and
 \(\overline P(f)=\max_{L\in\mathcal M}L(f)\), both attained, so
 \(\underline P(f)\leq\overline P(f)\);
\item[(P4)] \(\overline P(A)=1-\underline P(A^c)\) for \(A\in\Sigma\).
\end{enumerate}
Two consequences are used repeatedly. By (P1) and conjugacy,
\[
 \underline P(g-f)\geq\underline P(g)+\underline P(-f)
 =\underline P(g)-\overline P(f).\tag{M1}
\]
By (P1) applied to \(-g\) and \(g-f\), whose sum is \(-f\),
\(\underline P(-f)\geq\underline P(-g)+\underline P(g-f)\); negating,
\[
 \overline P(g)\geq\overline P(f)+\underline P(g-f).\tag{M2}
\]

\begin{annexitem}{annex:p15.w02b-p15-036.inclusions}{thm:p15-w02b-p15-036-inclusions}{The choice sets are nested}
\begin{proof}
\emph{Item 2.} Let \(f\in\operatorname{ch}_\alpha(\mathcal A)\) and
suppose some \(g\in\mathcal A\) has \(\delta=\underline P(g-f)>0\). By
(P1), \(\underline P(g)\geq\underline P(f)+\underline P(g-f)
=\underline P(f)+\delta\), and by (M2),
\(\overline P(g)\geq\overline P(f)+\delta\). Since \(\alpha\) and
\(1-\alpha\) are nonnegative and sum to one,
\(V_\alpha(g)\geq V_\alpha(f)+\delta>V_\alpha(f)\), contradicting the
choice of \(f\). So \(f\) is maximal.

\emph{Item 3.} Let \(f\in\operatorname{ch}_0(\mathcal A)\), so
\(\overline P(f)\geq\overline P(g)\) for every \(g\in\mathcal A\). By
(P3) there is \(L_f\in\mathcal M\) with \(L_f(f)=\overline P(f)\). For
every \(g\in\mathcal A\), again by (P3),
\(L_f(g)\leq\overline P(g)\leq\overline P(f)=L_f(f)\). So \(L_f\) ranks
\(f\) first on \(\mathcal A\), and \(f\) is E-admissible.

\emph{Item 4, first inclusion.} Let \(f\) be E-admissible with witness
\(L\in\mathcal M\). For every \(g\in\mathcal A\),
\(L(g-f)=L(g)-L(f)\leq0\) by linearity, so by (P3)
\(\underline P(g-f)\leq L(g-f)\leq0\). No \(g\) satisfies
\(\underline P(g-f)>0\), so \(f\) is maximal.

\emph{Item 4, second inclusion.} Let \(f\) be maximal and suppose some
\(g\in\mathcal A\) has \(\underline P(g)>\overline P(f)\). By (M1),
\(\underline P(g-f)\geq\underline P(g)-\overline P(f)>0\), contradicting
maximality. So \(f\) is interval undominated.

\emph{Item 1.} Each \(V_\alpha\) is a real function on the finite nonempty
set \(\mathcal A\), so it has a maximizer, and
\(\operatorname{ch}_\alpha(\mathcal A)\ne\varnothing\). By items~2 and~4,
\(\operatorname{ch}_1(\mathcal A)\subseteq\operatorname{ch}_{\mathrm{mx}}
(\mathcal A)\subseteq\operatorname{ch}_{\mathrm{id}}(\mathcal A)\), so the
last two sets are nonempty. Since \(\mathcal M\ne\varnothing\), pick
\(L\in\mathcal M\); the real function \(L\) has a maximizer on the finite
set \(\mathcal A\), and that act is E-admissible.

\emph{Item 5.} Let \(\mathcal A=\{f,g\}\) with \(f\ne g\) (if \(f=g\) the
menu has one act, which belongs to every set by item~1). Since
\(\underline P(f-f)=\underline P(0)=0\) by (P2), \(f\) is maximal exactly
when \(\underline P(g-f)\leq0\). It is E-admissible exactly when some
\(L\in\mathcal M\) has \(L(f-g)\geq0\). By (P3) the maximum of
\(L(f-g)\) over \(\mathcal M\) is attained and equals
\(\overline P(f-g)=-\underline P(g-f)\), so such an \(L\) exists exactly
when \(\underline P(g-f)\leq0\). The two conditions coincide for \(f\),
and by symmetry for \(g\). Both acts are therefore maximal exactly when
\(\underline P(g-f)\leq0\) and \(\underline P(f-g)\leq0\), that is, when
\(\underline P(f-g)\leq0\leq-\underline P(g-f)=\overline P(f-g)\).

\emph{Item 6.} Let \(\mathcal A=\{c_n:n\geq1\}\), where \(c_n\) is the
constant act \(1-1/n\); constants belong to \(B\). By (P2) and
conjugacy, \(\underline P(c_n)=\overline P(c_n)=c_n\), and
\(L(c_n)=c_n\) for every \(L\in\mathcal M\) because \(L\) is normalized and
linear. Hence \(V_\alpha(c_n)=c_n\), which has no maximum on
\(\mathcal A\), so \(\operatorname{ch}_\alpha(\mathcal A)=\varnothing\).
For every \(n\), \(\underline P(c_{n+1}-c_n)=1/n-1/(n+1)>0\), so no act is
maximal; \(L(c_{n+1})>L(c_n)\) for every \(L\), so no act is E-admissible;
and \(\underline P(c_{n+1})=c_{n+1}>c_n=\overline P(c_n)\), so no act is
interval undominated.
\end{proof}
\end{annexitem}

\begin{annexitem}{annex:p15.w02b-p15-036.precise}{thm:p15-w02b-p15-036-precise}{Agreement with expected utility exactly in the precise case}
\begin{proof}
\emph{(1) implies (2).} Let \(L\in\mathcal M\) and \(f\in B\). By
definition of the compatible set, \(L(f)\geq\underline P(f)\) and
\(L(-f)\geq\underline P(-f)=-\overline P(f)=-\underline P(f)\), so
\(L(f)\leq\underline P(f)\). Hence \(L(f)=\underline P(f)\) for every
\(f\), so every element of \(\mathcal M\) equals the functional
\(\underline P\). Since \(\mathcal M\ne\varnothing\), it has exactly one
element \(P\), and \(P=\underline P\) on \(B\).

\emph{(2) implies (1).} By (P3), \(\underline P(f)\) and
\(\overline P(f)\) are the minimum and the maximum of \(L(f)\) over the
one-element set \(\{P\}\), so both equal \(P(f)\).

\emph{(2) implies (3), with the stated common value.} Let
\(\mathcal M=\{P\}\) and let \(\mathcal A\) be a finite menu. By (1),
\(\underline P=\overline P=P\), so \(V_\alpha(f)=P(f)\) for every
\(\alpha\), and \(\operatorname{ch}_\alpha(\mathcal A)\) is the set of
maximizers of \(P\) on \(\mathcal A\). An act \(f\) is interval
undominated exactly when no \(g\in\mathcal A\) has \(P(g)>P(f)\); it is
maximal exactly when no \(g\) has \(P(g-f)=P(g)-P(f)>0\); and it is
E-admissible exactly when the only element \(P\) of \(\mathcal M\) has
\(P(f)\geq P(g)\) for every \(g\in\mathcal A\). Each of these is the
condition that \(f\) maximizes \(P\) on \(\mathcal A\). All members of
this set have the same value \(\max_{\mathcal A}P\).

\emph{(3) implies (1).} Suppose (1) fails: some \(f\in B\) has
\(\underline P(f)<\overline P(f)\). Then \(f\) is not constant, since a
constant has equal lower and upper previsions by (P2). Let \(c\) be the
constant act \((\underline P(f)+\overline P(f))/2\), so
\(\underline P(f)<c<\overline P(f)\) and \(f\ne c\). By (P2),
\(V_1(c)=\underline P(c)=c>\underline P(f)=V_1(f)\), so
\(\operatorname{ch}_1(\{f,c\})=\{c\}\). By (P2) and conjugacy,
\(\underline P(c-f)=c+\underline P(-f)=c-\overline P(f)<0\) and
\(\underline P(f-c)=\underline P(f)-c<0\), so neither act is strictly
preferred to the other and \(\operatorname{ch}_{\mathrm{mx}}(\{f,c\})=
\{f,c\}\). The two sets differ on this finite menu, so (3) fails.
\end{proof}
\end{annexitem}

\begin{annexitem}{annex:p15.w02b-p15-036.menu}{thm:p15-w02b-p15-036-menu}{The rules on the four-act menu}
\begin{proof}
\emph{Item 1.} Here \(B=\mathbb R^\Omega\), and a positive normalized
linear functional on it is a probability \(P_s\), \(s\in[0,1]\). Put
\(\underline P(h)=\min_{1/4\leq t\leq3/4}P_t(h)\); the minimum exists
because \(P_t(h)=t\,h(\omega_1)+(1-t)\,h(\omega_2)\) is affine in \(t\),
and it is attained at \(t=1/4\) or \(t=3/4\). The functional satisfies the
three conditions of \cref{thm:p15-ip-linear-domain}: each \(P_t(h)\geq
\inf h\), so the minimum is too; \(P_t(\lambda h)=\lambda P_t(h)\) for
\(\lambda>0\); and the minimum of \(P_t(h)+P_t(k)\) is at least the sum of
the separate minima. So \(\underline P\) is coherent. Each
\(P_t\in\mathcal M\) dominates \(\underline P\), so
\(\mathcal M\subseteq\mathcal M(\underline P)\). Conversely, if
\(P_s\in\mathcal M(\underline P)\) then
\(s=P_s(I_{\{\omega_1\}})\geq\underline P(I_{\{\omega_1\}})=1/4\) and
\(-s=P_s(-I_{\{\omega_1\}})\geq\underline P(-I_{\{\omega_1\}})
=\min_t(-t)=-3/4\), so \(s\in[1/4,3/4]\). Hence
\(\mathcal M(\underline P)=\mathcal M\).

The expectations are
\[
 P_t(a)=t,\qquad P_t(b)=\tfrac9{10}(1-t),\qquad P_t(c)=\tfrac25,\qquad
 P_t(d)=t-\tfrac1{20}.
\]
Each is affine in \(t\), so its minimum and maximum over \([1/4,3/4]\) are
its values at the endpoints: \(a\) gives \(1/4\) and \(3/4\); \(b\) gives
\(27/40\) at \(t=1/4\) and \(9/40\) at \(t=3/4\); \(c\) is constant
\(2/5\); \(d\) gives \(1/5\) and \(7/10\). This gives the four intervals.

\emph{Item 2.} From item~1,
\[
 V_\alpha(a)=\tfrac34-\tfrac\alpha2,\quad
 V_\alpha(b)=\tfrac{27}{40}-\tfrac{9\alpha}{20},\quad
 V_\alpha(c)=\tfrac25,\quad
 V_\alpha(d)=\tfrac7{10}-\tfrac\alpha2.
\]
Then \(V_\alpha(a)-V_\alpha(b)=(3-2\alpha)/40>0\) and
\(V_\alpha(a)-V_\alpha(d)=1/20>0\) for every \(\alpha\in[0,1]\), so
neither \(b\) nor \(d\) is ever chosen. Finally
\(V_\alpha(a)-V_\alpha(c)=7/20-\alpha/2\), which is positive for
\(\alpha<7/10\), zero at \(\alpha=7/10\), and negative for
\(\alpha>7/10\). This gives the three cases.

\emph{Item 3.} For every \(t\), \(P_t(a)-P_t(d)=1/20>0\), so \(d\) is best
under no \(P_t\). Next \(P_t(a)-P_t(b)=(19t-9)/10\), which is nonnegative
exactly when \(t\geq9/19\); note \(1/4<9/19<3/4\). So \(a\) is best among
\(\{a,b\}\) for \(t\in[9/19,3/4]\) and \(b\) for \(t\in[1/4,9/19]\). For
\(t\geq9/19\), \(P_t(a)=t\geq9/19\); for \(t<9/19\),
\(P_t(b)=\tfrac9{10}(1-t)>\tfrac9{10}\cdot\tfrac{10}{19}=9/19\). So under
every \(P_t\) the better of \(a\) and \(b\) has expectation at least
\(9/19\), which exceeds \(2/5=P_t(c)\) because \(45>38\). Hence \(c\) is
best under no \(P_t\). Combining, \(a\) is best on the whole menu for
\(t\in[9/19,3/4]\) and \(b\) for \(t\in[1/4,9/19]\), and by item~1 these
\(P_t\) are exactly the elements of \(\mathcal M\). So
\(\operatorname{ch}_{\mathrm E}(\mathcal A)=\{a,b\}\).

\emph{Item 4.} By item~2, \(c\in\operatorname{ch}_1(\mathcal A)\), and by
item~3, \(a,b\in\operatorname{ch}_{\mathrm E}(\mathcal A)\); by items~2
and~4 of \cref{thm:p15-w02b-p15-036-inclusions} all three are maximal.
The act \(d\) is not maximal, because \(a-d\) is the constant \(1/20\) and
\(\underline P(a-d)=1/20>0\) by (P2). So
\(\operatorname{ch}_{\mathrm{mx}}(\mathcal A)=\{a,b,c\}\). By the same
theorem these three acts are interval undominated. The largest lower
prevision on the menu is \(\underline P(c)=2/5\), and
\(\overline P(d)=7/10>2/5\), so no act interval-dominates \(d\). Hence
\(\operatorname{ch}_{\mathrm{id}}(\mathcal A)=\mathcal A\).

\emph{Item 5.} With all \(P_t\), \(0\leq t\leq1\), the lower envelope is
the map taking \(h\) to \(\min\{h(\omega_1),h(\omega_2)\}\); the argument of item~1, with
the interval \([0,1]\), shows it is coherent with compatible set
\(\{P_t:0\leq t\leq1\}\). Put \(u=(0,0)\) and \(v=(0,1)\). Then
\(P_t(v)-P_t(u)=1-t\geq0\) for every \(t\), with equality at \(t=1\) and
strict inequality at \(t=0\). So \(\underline P(v-u)=0\), which is not
positive, and \(\underline P(u-v)=-1\); both acts are maximal, and by
item~5 of \cref{thm:p15-w02b-p15-036-inclusions} both are E-admissible
(directly: \(P_1\) ranks \(u\) first, with a tie, and \(P_0\) ranks \(v\)
first). Under the strict part of unanimity, \(v\) is weakly preferred to
\(u\) under every \(P_t\), and the reverse weak preference fails under
\(P_0\), since \(P_0(u)=0<1=P_0(v)\). So \(v\) strictly dominates \(u\) in
that relation and \(u\) is excluded.
\end{proof}
\end{annexitem}

\begin{annexitem}{annex:p15.w02b-p15-036.accepted}{thm:p15-w02b-p15-036-accepted}{Full acceptance and deciding as if}
\begin{proof}
\emph{(1) implies (2).} Every \(L\in\mathcal M\) is positive and
normalized, so \(L(A)=L(I_A)\leq L(1)=1\). By (P3),
\(1=\underline P(A)=\min_{L\in\mathcal M}L(A)\), so \(L(A)\geq1\) for every
\(L\). Hence \(L(A)=1\).

\emph{(2) implies (3).} Let \(f,g\in B\) agree on \(A\) and put
\(h=f-g\). Then \(h=0\) on \(A\), so
\(-\|h\|_\infty I_{A^c}\leq h\leq\|h\|_\infty I_{A^c}\) pointwise. Since
\(A\in\Sigma\), \(I_{A^c}\in B\). For \(L\in\mathcal M\), positivity and
linearity give \(|L(h)|\leq\|h\|_\infty L(A^c)=\|h\|_\infty(1-L(A))=0\).
So \(L(f)=L(g)\) for every \(L\in\mathcal M\), and by (P3)
\(\underline P(f)=\underline P(g)\).

\emph{(3) implies (1).} The gambles \(I_A\) and \(1\) agree on \(A\), so
\(\underline P(A)=\underline P(1)=1\) by (P2).

\emph{The choice sets.} Assume (2) and let \(\varphi\) be as stated. By the
argument for (2) implies (3), \(L(\varphi(f))=L(f)\) for all
\(f\in\mathcal A\) and \(L\in\mathcal M\), and the same holds for
differences: \(L(\varphi(g)-\varphi(f))=L(g-f)\). Every quantity in
\cref{def:p15-w02b-p15-036-rules} is computed from the values \(L(f)\),
\(L\in\mathcal M\): by (P3) \(\underline P\) and \(\overline P\) of an act
or of a difference of acts are their minimum and maximum over
\(\mathcal M\), and E-admissibility compares \(L(f)\) with \(L(g)\).
Since \(\varphi\) is a bijection of \(\mathcal A\) onto
\(\varphi(\mathcal A)\), each quantifier over \(g\in\mathcal A\) in a
membership condition becomes the same quantifier over
\(\varphi(g)\in\varphi(\mathcal A)\), with the same values. Hence
\(f\) belongs to a given choice set for \(\mathcal A\) exactly when
\(\varphi(f)\) belongs to the corresponding set for \(\varphi(\mathcal A)\).

\emph{The failure case.} Suppose \(\underline P(A)<1\). Then
\(A^c\ne\varnothing\), so \(0\) and \(-I_{A^c}\) are distinct acts, and
they agree on \(A\). By conjugacy and (P4),
\(\underline P(-I_{A^c})=-\overline P(A^c)=-(1-\underline P(A))<0
=\underline P(0)\), using (P2) for the last value. So
\(V_1(0)>V_1(-I_{A^c})\) and
\(\operatorname{ch}_1(\{0,-I_{A^c}\})=\{0\}\).
\end{proof}
\end{annexitem}

\begin{annexitem}{annex:p15.w02b-p15-036.threshold}{thm:p15-w02b-p15-036-threshold}{A conjunction-closed acceptance that does not license deciding as if}
\begin{proof}
\emph{Item 1.} The argument of item~1 in the proof of
\cref{thm:p15-w02b-p15-036-menu}, with the interval \([9/10,1]\) in place
of \([1/4,3/4]\), shows that \(\underline P\) is coherent with compatible
set \(\mathcal M\). On events, \(\underline P(\varnothing)=0\),
\(\underline P(\Omega)=1\), \(\underline P(A)=\min_tt=9/10\) and
\(\underline P(A^c)=\min_t(1-t)=0\). The possibility measure of \(\pi\) is
\(\Pi(E)=\max_{\omega\in E}\pi(\omega)\), with empty maximum \(0\), and its
necessity measure is \(N(E)=1-\Pi(E^c)\). So \(N(\varnothing)=1-1=0\),
\(N(\Omega)=1-0=1\), \(N(A)=1-\pi(\omega_2)=9/10\) and
\(N(A^c)=1-\pi(\omega_1)=0\), which agree with \(\underline P\) on all
four events. The events of positive lower probability are \(A\) and
\(\Omega\); their intersections are \(A\) and \(\Omega\) again. And
\(\underline P(A)=9/10\geq\beta\) exactly when \(\beta\leq9/10\).

\emph{Item 2.} On \(A\), \(f(\omega_1)=1>0=g(\omega_1)\). For
\(t\in[9/10,1]\), \(P_t(f-g)=t-100(1-t)=101t-100\), which increases in
\(t\); its minimum, at \(t=9/10\), is \(909/10-100=-91/10\), and its
maximum, at \(t=1\), is \(1\). By (P3) these are \(\underline P(f-g)\) and
\(\overline P(f-g)\). Item~5 of \cref{thm:p15-w02b-p15-036-inclusions}
gives \(\operatorname{ch}_{\mathrm{mx}}(\{f,g\})=
\operatorname{ch}_{\mathrm E}(\{f,g\})=\{f,g\}\), since
\(-91/10\leq0\leq1\). Since \(g=0\), \(\underline P(f)=\underline P(f-g)
=-91/10<0=\underline P(g)\), so \(\operatorname{ch}_1(\{f,g\})=\{g\}\).

\emph{Item 3.} On the two-point space a probability with
\(P(A)=9/10\) is \(P_{9/10}\), and \(P_{9/10}(f)=-91/10<0=P_{9/10}(g)\) by
the computation in item~2.
\end{proof}
\end{annexitem}

\begin{annexitem}{annex:p15.w02b-p15-036.vacuous}{thm:p15-w02b-p15-036-vacuous}{The vacuous model and decision with only a set of possible states}
\begin{proof}
With \(\Omega\) finite and \(\Sigma=\mathcal P(\Omega)\), \(B=\mathbb
R^\Omega\). The functional \(\underline P(f)=\min f\) satisfies the three
conditions of \cref{thm:p15-ip-linear-domain}: it equals \(\inf f\), it is
positively homogeneous, and \(\min(f+g)\geq\min f+\min g\). So it is
coherent. A positive normalized linear functional \(L\) on
\(\mathbb R^\Omega\) has the form \(L(f)=\sum_\omega p(\omega)f(\omega)\)
with \(p(\omega)=L(I_{\{\omega\}})\geq0\) and \(\sum_\omega p(\omega)=
L(1)=1\), and every such \(L\) satisfies \(L(f)\geq\min f\). So
\(\mathcal M\) is the set of all mass functions. By conjugacy,
\(\overline P(f)=-\min(-f)=\max f\).

Item~1 follows by substituting these values in \(V_\alpha\). For item~2,
\(\underline P(g-f)=\min_\omega(g(\omega)-f(\omega))\), which is positive
exactly when \(g(\omega)>f(\omega)\) for every \(\omega\), the minimum over
the finite set \(\Omega\) being attained. Item~3 is
\cref{def:p15-w02b-p15-036-rules}, item~2, with
\(\underline P(g)=\min g\) and \(\overline P(f)=\max f\). Item~4 is
\cref{def:p15-w02b-p15-036-rules}, item~4, with \(\mathcal M\) the set of
all mass functions and \(L(f)=\sum_\omega p(\omega)f(\omega)\).
\end{proof}
\end{annexitem}
